\subsection{What the controls establish}
The added seeds reproduce the cohort-mean conditional preservation contrast; randomized maps show that equal RMS can support harmful as well as beneficial allocations; and the floor comparison links spatial redistribution to editing response. Together these results distinguish the allocation estimate from the choice of a complete editing configuration.
\paragraph{Seed replication on the same sources}
Averaging seeds 1, 2, and 3 gives REQ-U minus U25 preservation LPIPS $-0.01145$, CI $[-0.01714,-0.00592]$, Holm2 $p=0.000404$. REQ minus GLB-R gives $-0.01140$, CI $[-0.01572,-0.00735]$, Holm2 $p=0.000004$. Both profile-specific means are negative in each of the five observed seeds (Table~\ref{tab:r91-seeds}). Pooled five-seed means are $-0.01088$ and $-0.01082$, with marginal intervals $[-0.01582,-0.00618]$ and $[-0.01468,-0.00721]$. Both summaries reuse the same 40 sources; added seeds test sensitivity to generation noise, and five-seed pooling is descriptive.

The added-seed primary edit-MAE contrast is $+0.00808$, CI $[0.00592,0.01024]$, alongside a directional-CLIP contrast $+0.00075$, CI $[-0.00486,0.00572]$. Increased pixel change accompanies the preservation improvement again, while the directional-CLIP point difference changes sign relative to the original two seeds. The replicated joint observation is lower preservation error with greater edit-region departure; departure includes unintended changes, and directional-CLIP differences remain uncertain in both seed sets. Table~\ref{tab:r91-joint} reports the corresponding target CLIP and preservation MAE as well as both profile origins.

\paragraph{Individual-source variation}
For the primary contrast, 14 sources favor request weighting under all five seeds, one favors uniform under all five, and 25 have mixed signs. The five-seed source mean favors request weighting on 28/40 sources. Pairwise sign disagreements range from 9/40 to 17/40 across all ten seed pairs (Table~\ref{tab:r91-signs}); the original pair has 16/40. The five-seed mixed-sign count has more opportunities to change sign and is not a directly comparable instability rate. Figure~\ref{fig:r91-heterogeneity} shows every source/seed cell. The replicated cohort mean coexists with substantial source-specific variation.

\paragraph{Randomized placement and source-derived ranks}
Correct REQ improves over RAND by $0.01355$ preservation LPIPS: difference $-0.01355$, CI $[-0.01859,-0.00899]$, Holm3 $p=0.000003$. GLB-R also improves over RAND by $0.00359$ (CI $[-0.00514,-0.00217]$, Holm3 $p=0.000006$), and SOURCE-R improves by $0.00601$ (CI $[-0.00942,-0.00299]$, Holm3 $p=0.000636$). All three corresponding directional-CLIP intervals include zero (Table~\ref{tab:r91-permutation}). The correct-map versus RAND directional-CLIP difference is $-0.00133$, CI $[-0.00630,0.00336]$.

Relative to GLB-R, the SOURCE-R-to-REQ gain ratio remains 0.244, whereas the randomized-map ratio is $-0.360$: randomized weighting has higher preservation error than uniform correction at the same profile. SOURCE-R replaces spatial ranks with source-only velocity ranks while retaining the target-derived weight histogram and native request anchor. The ratios summarize mean outcome contrasts; 0.244 therefore measures the retained gain under that particular rank replacement, with no interpretation as a fraction or lower bound of source information.

The preserved histogram and matched RMS accommodate both beneficial and harmful spatial patterns. Cell permutation jointly disrupts smoothness and alignment with image structure, so the evidence concerns spatial organization as a whole. Highly correlated wrong-request maps yield uncertain contrasts against the correct request. Separating correct-request specificity from generic spatial structure remains open within these controls.

\paragraph{Floor sensitivity and joint outcomes}
On the fixed 20-source subset, floors 0, 0.25, and 0.5 give preservation LPIPS 0.46661, 0.46749, and 0.46888; its U25 reference gives 0.48026. Floor 0 minus U25 is $-0.0137$ LPIPS with descriptive 95\% interval $[-0.0251,-0.0023]$, and $+0.0033$ directional CLIP with descriptive interval $[-0.0033,0.0102]$. Floor 0.25 minus U25 directional CLIP is $-0.0051$ ($[-0.0121,0.0018]$). These paired comparisons use 20 sources, separate from the 40-source scale curves. Floor 0 has favorable joint point estimates; CLIP uncertainty leaves Pareto superiority unestablished.

Relative to floor 0.25, floor 0 changes LPIPS by $-0.00088$ (CI $[-0.00416,0.00242]$), and floor 0.5 by $+0.00139$ (CI $[-0.00155,0.00430]$); both four-test Holm-adjusted $p$ values are 1.0. Floor 0 raises directional CLIP by $+0.00850$ (CI $[0.00336,0.01407]$; exact four-test Holm-adjusted $p=0.0497$) and edit MAE by $+0.00724$ (CI $[0.00541,0.00888]$). The adjusted CLIP result is close to the 0.05 threshold. Floor 0.5 gives CLIP $-0.00129$ (CI $[-0.00567,0.00310]$; adjusted $p=1.0$) and edit MAE $-0.00423$ (CI $[-0.00545,-0.00295]$). LPIPS intervals spanning zero describe uncertainty, not equivalence.

Under normalized replay, the local multiplier is $b_i^U P_f(x)/\RMS(P_f\odot c_i)$; a weight floor of 0.25 is a relative spatial floor, not an absolute 25\% source correction. Floor 0 permits zero correction at maximum-response locations and changes the normalized direction while holding full-vector stepwise RMS fixed. Regional correction energy can change under this constraint. The result demonstrates floor-sensitive editing response at the tested profile; explaining the high-scale CLIP decline or estimating a floor-zero equal-preservation tradeoff would require further comparisons. Floor 0.25 remains the primary reference, and the fixed-subset outcomes are reported separately from the 40-source scale curves.
